\documentclass[12pt,a4paper]{article}
\usepackage[utf8]{inputenc}
\usepackage[T1]{fontenc}
\usepackage[spanish,es-tabla]{babel}
\usepackage{lmodern}
\usepackage[a4paper,margin=2.2cm]{geometry}
\usepackage{array}
\usepackage{longtable}
\usepackage{microtype}
\usepackage{hyperref}
\hypersetup{colorlinks=true,linkcolor=blue,urlcolor=blue}
\setlength{\parindent}{0pt}
\setlength{\parskip}{6pt}
\begin{document}
\begin{titlepage}
\centering
{\bfseries UNIVERSIDAD TÉCNICA ESTATAL DE QUEVEDO\par}
\vspace{0.7cm}
FACULTAD DE CIENCIAS DE LA COMPUTACIÓN\par
CARRERA DE SOFTWARE\par
INGENIERÍA DE REQUISITOS (ISR-401)\par
\vfill
{\Large\bfseries GA-PRAC-01 - FORM\par}
\vspace{0.5cm}
{\large Comparación entre modelos de proceso de requisitos en escenarios reales\par}
\vspace{1cm}
\textbf{Integrantes:}\par
Córdova Carreño Mayra Lucila\par
Herrera Morán Gilmar Javier\par
Rivas Bodero Jordy Fabian\par
\vspace{0.8cm}
\textbf{Documento integrado:} C1, C2, C3, C4 y C5\par
\vfill
Noviembre de 2026\par
\end{titlepage}
\tableofcontents
\clearpage

% GA-PRAC-01
% C1 - Caracterizacion de los modelos de proceso
% Responsable: Cordova Carreno Mayra Lucila

\section{C1. Caracterización de los modelos de proceso de requisitos}

\subsection{Modelo secuencial o cascada}

El modelo secuencial o cascada organiza el desarrollo del software mediante fases principalmente consecutivas. Las actividades de Ingeniería de Requisitos se realizan de manera planificada y sus resultados sirven como base para las etapas posteriores.

\textbf{Fases del desarrollo:} análisis de requisitos, diseño, implementación, pruebas, despliegue y mantenimiento. Las actividades de elicitación, análisis, especificación y validación permiten establecer los requisitos que orientan el desarrollo. La gestión de requisitos acompaña el proceso.

\textbf{Supuestos:} se considera que los requisitos pueden identificarse y documentarse con suficiente estabilidad durante las etapas iniciales. Los cambios posteriores deben gestionarse mediante procedimientos de control.

\textbf{Condiciones de aplicación:} resulta favorable cuando los requisitos son relativamente estables, existen obligaciones documentales, se necesitan aprobaciones formales y las actividades deben seguir una planificación definida.

\textbf{Limitaciones:} los cambios tardíos pueden provocar retrabajo y afectar los costos y plazos. Además, una validación funcional realizada demasiado tarde puede dificultar la detección oportuna de errores.

\subsection{Modelo iterativo e incremental}

El modelo iterativo e incremental organiza el desarrollo mediante ciclos sucesivos. En cada iteración se revisan y perfeccionan los requisitos, mientras que los incrementos permiten incorporar progresivamente funcionalidades al sistema.

\textbf{Fases o actividades:} planificación de la iteración, elicitación y análisis de requisitos, especificación, diseño, implementación, pruebas y evaluación de resultados. Estas actividades pueden repetirse durante diferentes ciclos.

\textbf{Supuestos:} los requisitos pueden comprenderse y perfeccionarse progresivamente mediante la evaluación de resultados parciales y la retroalimentación de las partes interesadas.

\textbf{Condiciones de aplicación:} resulta favorable cuando existe incertidumbre sobre algunos requisitos, se necesitan validaciones progresivas y es posible organizar el sistema en incrementos funcionales.

\textbf{Limitaciones:} requiere coordinación entre iteraciones, gestión de dependencias, control de cambios y planificación de las funcionalidades que se incorporarán en cada incremento.

\subsection{Modelo ágil}

El modelo ágil se caracteriza por su capacidad de adaptación, la colaboración frecuente con las partes interesadas y la entrega continua o frecuente de valor. Los requisitos se refinan y priorizan conforme evoluciona el conocimiento del problema.

\textbf{Fases o actividades:} identificación de necesidades, elicitación continua, priorización y refinamiento de requisitos, planificación de ciclos cortos, desarrollo, pruebas y validación frecuente. Estas actividades no necesariamente siguen una secuencia rígida.

\textbf{Supuestos:} los requisitos pueden cambiar durante el proyecto y existe disponibilidad suficiente de las partes interesadas para proporcionar retroalimentación y participar en las decisiones.

\textbf{Condiciones de aplicación:} resulta favorable cuando existe alta volatilidad de requisitos, se necesita responder rápidamente a cambios y es posible mantener una colaboración frecuente con los stakeholders.

\textbf{Limitaciones:} la participación insuficiente de los stakeholders, la falta de priorización y un control inadecuado del alcance pueden dificultar el cumplimiento de los objetivos. Los proyectos con exigencias contractuales o normativas requieren mecanismos adicionales de documentación y trazabilidad.

\subsection{Comparación general}

El modelo secuencial prioriza la planificación y la definición temprana de requisitos; el modelo iterativo e incremental permite perfeccionarlos y validar resultados progresivamente; mientras que el modelo ágil enfatiza la adaptación continua y la colaboración frecuente.

La elección del modelo debe considerar la estabilidad de los requisitos, la disponibilidad de las partes interesadas, la criticidad del sistema, los riesgos y las obligaciones documentales del proyecto.

\clearpage

% GA-PRAC-01
% C2 - Escenario 1: Extractora Rio Lindo
% Responsable: Cordova Carreno Mayra Lucila
% Requiere: \usepackage{longtable}
%           \usepackage{array}

\section{C2. Matriz comparativa contextualizada}

\subsection{Escenario 1: Extractora Río Lindo}

La Extractora Río Lindo requiere un sistema embebido
para controlar temperatura, presión y niveles de
tanques. El proyecto presenta requisitos de seguridad
estables e incertidumbre en sensores, interfaz y
reportes. Su criticidad es muy alta debido a los
riesgos para los operadores y la producción.

\begingroup
\footnotesize
\setlength{\tabcolsep}{4pt}
\renewcommand{\arraystretch}{1.24}
\setlength{\LTleft}{0pt}
\setlength{\LTright}{0pt}

\begin{longtable}{|>{\raggedright\arraybackslash}p{0.155\textwidth}|>{\raggedright\arraybackslash}p{0.255\textwidth}|>{\raggedright\arraybackslash}p{0.255\textwidth}|>{\raggedright\arraybackslash}p{0.255\textwidth}|}
\hline
\textbf{Criterio} &
\textbf{Secuencial (cascada)} &
\textbf{Iterativo e incremental} &
\textbf{Ágil} \\
\hline
\endfirsthead

\hline
\textbf{Criterio} &
\textbf{Secuencial (cascada)} &
\textbf{Iterativo e incremental} &
\textbf{Ágil} \\
\hline
\endhead

\textbf{K1. Estabilidad y cambios} &
Favorece la especificación temprana de temperaturas,
presiones y bloqueos, pero dificulta incorporar
cambios posteriores en sensores e interfaz. &
Permite mantener estables los requisitos de seguridad
y refinar por iteraciones la comunicación con sensores
y los reportes. &
Facilita adaptar interfaz y reportes, pero los cambios
frecuentes no aportan la misma ventaja a las reglas
de seguridad ya definidas. \\
\hline

\textbf{K2. Partes interesadas y disponibilidad} &
Permite planificar reuniones con el jefe de planta
dos mañanas por semana y obtener aprobaciones
formales del responsable de seguridad. &
Permite programar revisiones con el jefe de planta
y operadores de distintos turnos al finalizar
cada iteración. &
La colaboración frecuente resulta difícil por la
disponibilidad limitada del jefe de planta y los
turnos rotativos de los operadores. \\
\hline

\textbf{K3. Tipo de sistema y criticidad} &
Favorece la trazabilidad de los requisitos de
protección del sistema embebido, aunque puede
retrasar pruebas de integración. &
Permite comprobar progresivamente alarmas, bloqueos
y comunicación con dispositivos, conservando
requisitos de seguridad verificables. &
Facilita ajustar la interfaz, pero necesita controles
rigurosos para evitar que cambios rápidos comprometan
funciones críticas de protección. \\
\hline

\textbf{K4. Riesgo principal} &
La selección de sensores en el tercer mes podría
invalidar decisiones técnicas anteriores y
generar retrabajo. &
Permite verificar tempranamente la integración
con sensores y válvulas, reduciendo incertidumbre
antes de la puesta en marcha. &
Existe riesgo de priorizar mejoras visibles de
interfaz sobre verificaciones críticas si no se
establecen prioridades y criterios de seguridad
estrictos. \\
\hline

\textbf{K5. Documentación y marco normativo o contractual} &
Favorece la aprobación formal de requisitos de
seguridad y la documentación de pruebas exigidas
por la aseguradora. &
Permite conservar versiones de requisitos y
evidencias de pruebas por incremento, manteniendo
la trazabilidad de funciones de protección. &
Requiere complementar la documentación habitual
de trabajo con requisitos de seguridad aprobados
y evidencias verificables de cada función de
protección. \\
\hline

\end{longtable}
\endgroup

\normalsize

\clearpage

% GA-PRAC-01
% C2 - Escenario 2: Sabores de la Abuela Rosa
% Responsable: Cordova Carreno Mayra Lucila
% Requiere: longtable y array

\subsection{Escenario 2: Sabores de la Abuela Rosa}

El emprendimiento Sabores de la Abuela Rosa necesita
una aplicación móvil para gestionar pedidos de comida
típica. El proyecto presenta alta volatilidad de
requisitos, disponibilidad frecuente de la propietaria,
un presupuesto pequeño y un plazo de 14 semanas.

El sistema debe proteger los datos personales de los
clientes, controlar los pagos por transferencia y
funcionar en teléfonos Android de gama baja con
conectividad intermitente.

\begingroup
\footnotesize
\setlength{\tabcolsep}{4pt}
\renewcommand{\arraystretch}{1.24}
\setlength{\LTleft}{0pt}
\setlength{\LTright}{0pt}

\begin{longtable}{|>{\raggedright\arraybackslash}p{0.155\textwidth}|>{\raggedright\arraybackslash}p{0.255\textwidth}|>{\raggedright\arraybackslash}p{0.255\textwidth}|>{\raggedright\arraybackslash}p{0.255\textwidth}|}
\hline
\textbf{Criterio} &
\textbf{Secuencial (cascada)} &
\textbf{Iterativo e incremental} &
\textbf{Ágil} \\
\hline
\endfirsthead

\hline
\textbf{Criterio} &
\textbf{Secuencial (cascada)} &
\textbf{Iterativo e incremental} &
\textbf{Ágil} \\
\hline
\endhead

\textbf{K1. Estabilidad y cambios} &
Una especificación inicial quedaría desactualizada
por los cambios diarios del menú y las nuevas ideas
de Doña Rosa y su hijo. &
Permite incorporar progresivamente cambios en
promociones, menú y funcionalidades mediante
ciclos de revisión. &
Favorece la repriorización frecuente de requisitos
según los cambios del menú y las decisiones
de Doña Rosa. \\
\hline

\textbf{K2. Partes interesadas y disponibilidad} &
Doña Rosa podría aprobar requisitos al inicio,
pero su disponibilidad diaria quedaría poco
aprovechada durante el desarrollo. &
Permite que Doña Rosa pruebe los incrementos
y comunique correcciones al finalizar cada ciclo. &
Aprovecha su disponibilidad diaria para validar
funcionalidades y decidir prioridades, aunque
deben controlarse las solicitudes adicionales
de su hijo. \\
\hline

\textbf{K3. Tipo de sistema y criticidad} &
Permite documentar requisitos de pagos y protección
de datos, pero podría detectar tarde problemas
de uso en teléfonos económicos. &
Facilita probar progresivamente pedidos, pagos
y funcionamiento con conexiones intermitentes. &
Permite validar tempranamente la experiencia móvil,
pero exige criterios de aceptación para pagos,
privacidad y funcionamiento sin conexión estable. \\
\hline

\textbf{K4. Riesgo principal} &
Los cambios frecuentes podrían provocar retrabajo
y comprometer la entrega dentro de las 14 semanas. &
Existe riesgo de no completar los incrementos
previstos con dos desarrolladores a medio tiempo. &
Las nuevas ideas pueden ampliar excesivamente
el alcance y agotar el presupuesto antes
de las fiestas. \\
\hline

\textbf{K5. Documentación y marco normativo o contractual} &
Favorece una especificación formal de datos
personales y pagos, aunque requiere actualizarla
ante cambios frecuentes. &
Permite mantener requisitos y evidencias
de pruebas actualizados por incremento. &
Puede utilizar historias de usuario y criterios
de aceptación, complementados con documentación
de privacidad y conciliación de transferencias. \\
\hline

\end{longtable}
\endgroup

\normalsize

\clearpage

% GA-PRAC-01
% C2 - Escenario 3: Sistema de tramites municipales
% Responsable: Cordova Carreno Mayra Lucila
% Requiere: longtable y array

\subsection{Escenario 3: Sistema de trámites municipales contratado por contratación pública}

El GAD Municipal del cantón San Isidro de Baba, en la provincia de Los Ríos, requiere desarrollar un sistema web para gestionar aproximadamente 9\,000 trámites anuales relacionados con permisos de funcionamiento, patentes y certificados de no adeudar.

El proyecto será ejecutado por una empresa contratista de seis personas, bajo un contrato público de siete meses. Se establecen tres entregables con fechas y pagos fijos: especificación de requisitos aprobada en el mes 2, sistema instalado y probado en el mes 6, y capacitación con acta de entrega-recepción definitiva en el mes 7. Cada día de retraso genera una multa.

Los requisitos presentan incertidumbre debido a posibles cambios en el flujo de aprobación y a una nueva ordenanza de tasas que podría aprobarse en el cuarto mes. Además, el sistema debe proteger los datos personales, administrar cobros y conectarse con servidores municipales y una base de datos antigua sin documentación.

\begingroup
\footnotesize
\setlength{\tabcolsep}{4pt}
\renewcommand{\arraystretch}{1.24}
\setlength{\LTleft}{0pt}
\setlength{\LTright}{0pt}

\begin{longtable}{|>{\raggedright\arraybackslash}p{0.155\textwidth}|>{\raggedright\arraybackslash}p{0.255\textwidth}|>{\raggedright\arraybackslash}p{0.255\textwidth}|>{\raggedright\arraybackslash}p{0.255\textwidth}|}
\hline
\textbf{Criterio} &
\textbf{Cascada} &
\textbf{Iterativo e incremental} &
\textbf{Ágil} \\
\hline
\endfirsthead

\hline
\textbf{Criterio} &
\textbf{Cascada} &
\textbf{Iterativo e incremental} &
\textbf{Ágil} \\
\hline
\endhead

\textbf{K1. Estabilidad y cambios} &
Permite establecer una línea base de requisitos para la aprobación contractual del mes 2. Sin embargo, los cambios en el flujo de aprobación o en las tasas municipales requerirían una evaluación formal de impacto y posiblemente una modificación contractual. &
Permite revisar y ajustar los requisitos mediante ciclos de desarrollo, incorporando los cambios previstos para el cuarto mes. No obstante, los ajustes deben controlarse para respetar los entregables y pagos establecidos. &
Facilita adaptar los requisitos mediante la revisión continua del backlog. Sin embargo, los cambios solicitados por los funcionarios no pueden incorporarse libremente si modifican el alcance contractual aprobado. \\
\hline

\textbf{K2. Partes interesadas y disponibilidad} &
La participación del administrador del contrato y de los funcionarios se concentra principalmente en la especificación, revisión y aceptación formal. La disponibilidad semanal del administrador favorece reuniones programadas, pero limita la retroalimentación continua. &
Permite realizar validaciones periódicas con el administrador y los funcionarios de ventanilla y rentas. Las revisiones pueden organizarse según su disponibilidad y los hitos del contrato. &
Requiere colaboración frecuente con representantes del municipio. La participación irregular de los funcionarios y la disponibilidad limitada del administrador dificultan mantener una retroalimentación ágil constante. \\
\hline

\textbf{K3. Tipo de sistema y criticidad} &
La secuencia de fases facilita documentar y verificar las funciones de permisos, patentes y certificados. Sin embargo, descubrir tardíamente problemas de integración con la base de datos antigua podría afectar el cronograma. &
Permite desarrollar y validar progresivamente los módulos de trámites, cobros e integración. Las primeras iteraciones pueden utilizarse para investigar la base de datos antigua y comprobar la compatibilidad con los servidores municipales. &
Permite entregar funcionalidades parciales y recibir observaciones tempranas. No obstante, el sistema necesita cumplir requisitos de seguridad, integración y aceptación contractual que exigen controles adicionales. \\
\hline

\textbf{K4. Riesgo principal} &
Presenta un riesgo elevado de retrabajo si los cambios normativos aparecen después de aprobar la especificación. Los problemas técnicos descubiertos durante la implementación podrían ocasionar multas por retraso. &
Permite identificar anticipadamente riesgos de integración, protección de datos y cálculos de tasas. Sin embargo, requiere planificar cuidadosamente las iteraciones para evitar incumplir los tres hitos contractuales. &
Permite detectar problemas mediante pruebas y revisiones frecuentes. Sin embargo, una priorización inadecuada del backlog o la falta de disponibilidad del cliente podría comprometer las fechas y el alcance contratado. \\
\hline

\textbf{K5. Documentación y marco normativo o contractual} &
Favorece una especificación formal, trazabilidad de requisitos, registros de pruebas y actas de aceptación, necesarios para justificar los entregables y pagos del contrato público. &
Permite mantener una especificación inicial aprobada y actualizar los artefactos mediante control de versiones, trazabilidad y registros de cambios, conservando la evidencia exigida por el contrato. &
Utiliza historias de usuario, criterios de aceptación y un backlog priorizado. Para este contrato sería necesario complementar esos artefactos con especificaciones formales, trazabilidad y documentación de aceptación. \\
\hline

\end{longtable}
\endgroup

\normalsize

\noindent\textbf{Síntesis comparativa.}

El modelo cascada ofrece mayor correspondencia con la aprobación formal de requisitos y los entregables contractuales, pero presenta dificultades ante los cambios normativos y la incertidumbre técnica.

El modelo iterativo e incremental permite reducir progresivamente los riesgos de integración y validar los requisitos con los funcionarios municipales, siempre que sus ciclos se ajusten a los hitos y procedimientos formales del contrato.

El modelo ágil ofrece flexibilidad para responder a cambios, pero la disponibilidad limitada de los stakeholders y las restricciones de contratación pública dificultan su aplicación sin mecanismos adicionales de control documental y contractual.

\clearpage

% GA-PRAC-01
% C2 - Escenario 4: Cooperativa Frutos del Estero
% Responsable: Cordova Carreno Mayra Lucila
% Requiere: longtable y array

\subsection{Escenario 4: Plataforma web para una cooperativa de productores de cacao con socios dispersos}

La Cooperativa Agroproductiva Frutos del Estero, ubicada en la provincia de Los Ríos, agrupa a 640 pequeños productores de cacao de recintos rurales de Mocache, Ventanas y Valencia.

La cooperativa necesita una plataforma web que permita registrar las entregas de cacao según peso, humedad y calidad, calcular las liquidaciones y consultar entregas, pagos, créditos y saldos. Actualmente, estos procesos se realizan mediante tarjetas de papel y hojas de cálculo.

El proyecto dispone de un equipo de cuatro personas y un plazo de doce meses. La gerente está disponible casi diariamente, pero los cambios importantes requieren aprobación de la directiva, que se reúne mensualmente. Los socios se encuentran dispersos y participan principalmente durante los días de acopio o las asambleas trimestrales.

Las reglas de liquidación presentan baja volatilidad, mientras que los requisitos de trazabilidad y certificación del cacao pueden cambiar según las exigencias de los compradores exportadores. Además, el sistema debe proteger los datos personales, funcionar sin conexión durante el registro de entregas y sincronizar la información posteriormente.

El proyecto cuenta con financiamiento de un donante que exige informes semestrales de avance y está sujeto al estatuto de la cooperativa y al control del organismo supervisor de la economía popular y solidaria.

\begingroup
\footnotesize
\setlength{\tabcolsep}{4pt}
\renewcommand{\arraystretch}{1.24}
\setlength{\LTleft}{0pt}
\setlength{\LTright}{0pt}

\begin{longtable}{|>{\raggedright\arraybackslash}p{0.155\textwidth}|>{\raggedright\arraybackslash}p{0.255\textwidth}|>{\raggedright\arraybackslash}p{0.255\textwidth}|>{\raggedright\arraybackslash}p{0.255\textwidth}|}

\hline
\textbf{Criterio} &
\textbf{Cascada} &
\textbf{Iterativo e incremental} &
\textbf{Ágil} \\
\hline
\endfirsthead

\hline
\textbf{Criterio} &
\textbf{Cascada} &
\textbf{Iterativo e incremental} &
\textbf{Ágil} \\
\hline
\endhead

\textbf{K1. Estabilidad y cambios} &
Permite definir y documentar inicialmente las reglas de liquidación, que presentan baja volatilidad. Sin embargo, los cambios en trazabilidad y certificación solicitados por compradores exportadores podrían generar retrabajo y modificaciones formales de requisitos. &
Permite establecer una línea base para las reglas de liquidación y revisar progresivamente los requisitos de trazabilidad y certificación. Los cambios pueden incorporarse en nuevas iteraciones, considerando su impacto en el plazo de doce meses. &
Facilita actualizar y priorizar los requisitos de trazabilidad mediante un backlog. No obstante, los cambios importantes deben esperar la aprobación mensual de la directiva, lo que limita la capacidad de respuesta inmediata. \\
\hline

\textbf{K2. Partes interesadas y disponibilidad} &
Permite recopilar requisitos mediante entrevistas y reuniones iniciales con la gerente, la directiva y representantes de los socios. Sin embargo, la participación limitada de los productores podría ocasionar que sus necesidades reales se descubran durante las pruebas finales. &
Permite validar periódicamente las funcionalidades con la gerente y aprovechar los días de acopio y las asambleas trimestrales para obtener retroalimentación de los socios. Las decisiones importantes se presentan a la directiva para su aprobación. &
La disponibilidad casi diaria de la gerente favorece la colaboración continua. Sin embargo, ella no puede aprobar cambios importantes por sí sola y los socios dispersos no pueden participar frecuentemente en revisiones, lo que dificulta una validación ágil representativa. \\
\hline

\textbf{K3. Tipo de sistema y criticidad} &
Facilita especificar detalladamente los cálculos de liquidaciones, créditos y saldos antes de la implementación. Sin embargo, desarrollar la plataforma mediante fases principalmente consecutivas puede retrasar la validación del funcionamiento sin conexión y la sincronización de registros. &
Permite construir progresivamente los módulos de entregas, liquidaciones, créditos, consultas y trazabilidad. Las primeras iteraciones pueden comprobar el registro sin conexión y la sincronización posterior, reduciendo la incertidumbre técnica. &
Permite desarrollar funcionalidades prioritarias y demostrar avances frecuentes. Puede comenzar con el registro de entregas y continuar con pagos, créditos y trazabilidad, aunque requiere coordinar las dependencias entre módulos y las validaciones de los socios. \\
\hline

\textbf{K4. Riesgo principal} &
Presenta riesgo de detectar tardíamente errores en liquidaciones o problemas de sincronización. Estos fallos podrían afectar los pagos a los productores y la integridad de los registros de entregas. &
Permite identificar y corregir tempranamente errores de cálculo, duplicación de registros, pérdida de información y problemas de conectividad. Sin embargo, exige controlar las dependencias técnicas entre incrementos y mantener pruebas de integración. &
Las pruebas frecuentes permiten detectar problemas de usabilidad y funcionamiento. Sin embargo, un backlog mal priorizado podría postergar controles críticos de liquidación, seguridad de datos o sincronización, especialmente con un equipo de cuatro personas. \\
\hline

\textbf{K5. Documentación y marco normativo o contractual} &
Favorece una especificación formal de requisitos, reglas de negocio, documentación de cálculos, pruebas y evidencias para los informes semestrales exigidos por el donante y el control institucional. &
Permite mantener requisitos documentados y trazables, actualizar especificaciones por iteración y generar evidencias de pruebas y avances. Los resultados de cada incremento pueden respaldar los informes semestrales del donante. &
Utiliza historias de usuario, criterios de aceptación y backlog priorizado. Para cumplir las condiciones del donante, el estatuto y el control institucional, sería necesario complementar estos artefactos con documentación formal, trazabilidad y evidencias verificables. \\
\hline

\end{longtable}
\endgroup

\normalsize

\noindent\textbf{Síntesis comparativa.}

El modelo cascada ofrece ventajas para documentar las reglas estables de liquidación y cumplir las obligaciones institucionales. Sin embargo, su menor flexibilidad frente a cambios en trazabilidad y certificación, junto con la validación tardía de la sincronización, representa una limitación importante.

El modelo iterativo e incremental permite desarrollar y validar progresivamente los módulos de la plataforma, comprobar anticipadamente el funcionamiento sin conexión y adaptar los requisitos de trazabilidad. Además, facilita organizar revisiones con los socios según su disponibilidad y conservar evidencias para los informes semestrales.

El modelo ágil favorece la adaptación a cambios y la colaboración frecuente con la gerente. No obstante, la aprobación mensual de la directiva, la dispersión de los socios y las exigencias documentales del donante requieren mecanismos adicionales de coordinación y control.

\clearpage



% GA-PRAC-01
% C3 - Actores del proceso y framework de Ingenieria de Requisitos
% Responsable: Herrera Moran Gilmar Javier
% Requiere: \usepackage{longtable}
%           \usepackage{array}

\section{C3. Actores del proceso y framework de Ingeniería de Requisitos}

\subsection{Marco de referencia}

El framework de Ingeniería de Requisitos (IR) \cite{Pohl2025,IREB2026} se organiza en tres grupos de elementos que se aplican a los cuatro escenarios, independientemente del modelo de proceso elegido:

\begin{itemize}
    \item \textbf{Actividades core:} elicitación, documentación y negociación o resolución de conflictos. Producen y consolidan los requisitos \cite{Pohl2025}.
    \item \textbf{Actividades transversales:} validación y gestión de requisitos. Acompañan a las tres actividades core durante todo el proyecto \cite{Pohl2025,SWEBOK_v4_2024}.
    \item \textbf{Artefactos:} metas, escenarios y requisitos orientados a solución, que representan el problema y la solución con distinto nivel de detalle \cite{Pohl2025}.
\end{itemize}

En cada escenario se identifican primero los \textbf{actores del proceso}, es decir, quienes aportan, deciden, analizan, validan o implementan requisitos. Luego se ubican las cinco actividades del framework y, finalmente, se indica cómo se concretan los tres artefactos. La identificación de actores como partes interesadas y la trazabilidad de los requisitos siguen además las prácticas de la norma ISO/IEC/IEEE 29148 \cite{ISO29148}. Las tablas siguientes usan únicamente los datos de cada escenario y los modelos recomendados en C4.

% ---------------------------------------------------------------------
% ESCENARIO 1
% ---------------------------------------------------------------------
\subsection{Escenario 1: Extractora Río Lindo}

\textbf{Actores y responsabilidades.}

{\small
\setlength{\tabcolsep}{4pt}
\renewcommand{\arraystretch}{1.3}
\begin{longtable}{
|>{\raggedright\arraybackslash}p{0.22\textwidth}
|>{\raggedright\arraybackslash}p{0.16\textwidth}
|>{\raggedright\arraybackslash}p{0.50\textwidth}|}
\hline
\textbf{Actor} & \textbf{Tipo} & \textbf{Responsabilidad en el proceso de IR} \\
\hline
\endfirsthead
\hline
\textbf{Actor} & \textbf{Tipo} & \textbf{Responsabilidad en el proceso de IR} \\
\hline
\endhead
Jefe de planta & Stakeholder / decisor & Aporta las reglas operativas del proceso, aprueba requisitos y prioriza. Está disponible solo dos mañanas por semana, por lo que sus sesiones deben planificarse. \\
\hline
Responsable de seguridad & Stakeholder / aprobador & Define y aprueba los requisitos de protección (alarmas, bloqueos, límites de temperatura y presión). \\
\hline
Operadores de turnos rotativos & Usuarios finales & Describen el uso real de la interfaz y de las alarmas. Validan escenarios de operación en cada turno. \\
\hline
Aseguradora & Stakeholder externo & Exige un documento de requisitos de seguridad aprobado y evidencias de pruebas antes de la puesta en marcha. \\
\hline
Fabricante de los esterilizadores & Fuente de requisitos & Establece los límites técnicos y de seguridad del equipo que el sistema debe respetar. \\
\hline
Proveedores de sensores y válvulas & Stakeholder técnico & Su selección (tercer mes) determina interfaces y restricciones de integración. \\
\hline
Analista de requisitos & Equipo del proyecto & Conduce la elicitación, redacta la especificación, gestiona conflictos y mantiene la trazabilidad. \\
\hline
Equipo de desarrollo e integración & Equipo del proyecto & Implementa el sistema embebido, estima la factibilidad y verifica los requisitos en cada incremento. \\
\hline
\end{longtable}
}

\textbf{Actividades del framework ubicadas en el escenario.}

{\small
\setlength{\tabcolsep}{4pt}
\renewcommand{\arraystretch}{1.3}
\begin{longtable}{
|>{\raggedright\arraybackslash}p{0.20\textwidth}
|>{\raggedright\arraybackslash}p{0.13\textwidth}
|>{\raggedright\arraybackslash}p{0.55\textwidth}|}
\hline
\textbf{Actividad} & \textbf{Tipo} & \textbf{Aplicación en el escenario} \\
\hline
\endfirsthead
\hline
\textbf{Actividad} & \textbf{Tipo} & \textbf{Aplicación en el escenario} \\
\hline
\endhead
Elicitación & Core & Entrevistas con el jefe de planta en sus dos mañanas semanales, observación de operadores en distintos turnos y análisis de las especificaciones del fabricante y de la aseguradora. \\
\hline
Documentación & Core & Especificación formal de los requisitos de seguridad (umbrales, bloqueos y alarmas), seguida de la documentación incremental de interfaz, reportes y comunicación con sensores. \\
\hline
Negociación y resolución de conflictos & Core & Equilibrar el deseo de operadores de una interfaz cómoda con las restricciones de seguridad, y los plazos del proyecto con la incertidumbre de los sensores. \\
\hline
Validación & Transversal & Aprobación formal de los requisitos de seguridad por el responsable de seguridad y la aseguradora. Revisión de prototipos de interfaz con operadores de varios turnos al cierre de cada incremento. \\
\hline
Gestión de requisitos & Transversal & Línea base congelada para las reglas de seguridad, trazabilidad entre requisito, prueba y evidencia, y control formal de cambios cuando se definan los sensores. \\
\hline
\end{longtable}
}

\textbf{Artefactos del framework.}

{\small
\setlength{\tabcolsep}{4pt}
\renewcommand{\arraystretch}{1.3}
\begin{longtable}{
|>{\raggedright\arraybackslash}p{0.20\textwidth}
|>{\raggedright\arraybackslash}p{0.68\textwidth}|}
\hline
\textbf{Artefacto} & \textbf{Contenido en el escenario} \\
\hline
\endfirsthead
\hline
\textbf{Artefacto} & \textbf{Contenido en el escenario} \\
\hline
\endhead
Metas & Proteger a los operadores y la producción; controlar temperatura, presión y niveles de tanques dentro de los límites del fabricante; cumplir las exigencias de la aseguradora. \\
\hline
Escenarios & Un operador recibe una alarma por presión elevada y el sistema activa el bloqueo. Un operador de otro turno consulta el reporte del lote. Se pierde la lectura de un sensor y el sistema pasa a estado seguro. \\
\hline
Requisitos orientados a solución & Requisitos funcionales de lectura, alarma y bloqueo; requisitos de calidad como tiempo de respuesta y disponibilidad; restricciones de integración con los sensores y válvulas que se seleccionen. \\
\hline
\end{longtable}
}

% ---------------------------------------------------------------------
% ESCENARIO 2
% ---------------------------------------------------------------------
\subsection{Escenario 2: Sabores de la Abuela Rosa}

\textbf{Actores y responsabilidades.}

{\small
\setlength{\tabcolsep}{4pt}
\renewcommand{\arraystretch}{1.3}
\begin{longtable}{
|>{\raggedright\arraybackslash}p{0.22\textwidth}
|>{\raggedright\arraybackslash}p{0.16\textwidth}
|>{\raggedright\arraybackslash}p{0.50\textwidth}|}
\hline
\textbf{Actor} & \textbf{Tipo} & \textbf{Responsabilidad en el proceso de IR} \\
\hline
\endfirsthead
\hline
\textbf{Actor} & \textbf{Tipo} & \textbf{Responsabilidad en el proceso de IR} \\
\hline
\endhead
Doña Rosa (propietaria) & Cliente / decisora & Aporta las necesidades del negocio, prioriza el contenido de cada ciclo y valida los incrementos. Está disponible a diario. \\
\hline
Hijo de la propietaria & Stakeholder & Propone ideas y mejoras. Sus solicitudes deben pasar por la priorización de Doña Rosa para evitar ampliar el alcance. \\
\hline
Clientes del emprendimiento & Usuarios finales & Realizan pedidos y pagos por transferencia desde teléfonos Android de gama baja. Sus datos personales deben protegerse. \\
\hline
Personal de cocina o entrega & Usuarios internos & Reciben y actualizan el estado de los pedidos. \\
\hline
Dos desarrolladores a medio tiempo & Equipo del proyecto & Asumen elicitación, análisis, implementación y pruebas, con capacidad limitada. Uno de ellos actúa como analista de requisitos. \\
\hline
Autoridad de protección de datos (LOPDP) & Marco normativo & Establece las obligaciones sobre el tratamiento de datos personales de los clientes. \\
\hline
\end{longtable}
}

\textbf{Actividades del framework ubicadas en el escenario.}

{\small
\setlength{\tabcolsep}{4pt}
\renewcommand{\arraystretch}{1.3}
\begin{longtable}{
|>{\raggedright\arraybackslash}p{0.20\textwidth}
|>{\raggedright\arraybackslash}p{0.13\textwidth}
|>{\raggedright\arraybackslash}p{0.55\textwidth}|}
\hline
\textbf{Actividad} & \textbf{Tipo} & \textbf{Aplicación en el escenario} \\
\hline
\endfirsthead
\hline
\textbf{Actividad} & \textbf{Tipo} & \textbf{Aplicación en el escenario} \\
\hline
\endhead
Elicitación & Core & Conversaciones diarias breves con Doña Rosa, observación de cómo toma pedidos hoy, y consulta a algunos clientes sobre cómo piden y pagan. \\
\hline
Documentación & Core & Historias de usuario con criterios de aceptación en un \textit{backlog} vivo. Documentación breve y obligatoria de privacidad de datos y de conciliación de transferencias. \\
\hline
Negociación y resolución de conflictos & Core & Decidir qué ideas del hijo entran en cada ciclo y cuáles se aplazan, según presupuesto y plazo de 14 semanas. Resolver el choque entre agilidad comercial y cumplimiento legal. \\
\hline
Validación & Transversal & Doña Rosa prueba cada incremento en teléfonos de gama baja y con conexión intermitente. Se verifica la protección de datos y el control de pagos antes de cerrar cada ciclo. \\
\hline
Gestión de requisitos & Transversal & Priorización continua (por ejemplo, MoSCoW), control del alcance de cada ciclo y registro de cambios del menú y de las promociones. \\
\hline
\end{longtable}
}

\textbf{Artefactos del framework.}

{\small
\setlength{\tabcolsep}{4pt}
\renewcommand{\arraystretch}{1.3}
\begin{longtable}{
|>{\raggedright\arraybackslash}p{0.20\textwidth}
|>{\raggedright\arraybackslash}p{0.68\textwidth}|}
\hline
\textbf{Artefacto} & \textbf{Contenido en el escenario} \\
\hline
\endfirsthead
\hline
\textbf{Artefacto} & \textbf{Contenido en el escenario} \\
\hline
\endhead
Metas & Recibir y gestionar pedidos de comida típica de forma ordenada; controlar los pagos por transferencia; proteger los datos de los clientes; tener la aplicación lista para las fiestas dentro de 14 semanas y del presupuesto. \\
\hline
Escenarios & Un cliente elige un plato del menú del día y confirma el pedido. Un cliente con mala conexión intenta pagar y el sistema conserva el pedido. Doña Rosa cambia el menú y las promociones. \\
\hline
Requisitos orientados a solución & Funcionales: catálogo editable, carrito, registro de pedidos y de transferencias. De calidad: funcionamiento en Android de gama baja y con conectividad intermitente. Restricciones: protección de datos personales y presupuesto cerrado. \\
\hline
\end{longtable}
}

% ---------------------------------------------------------------------
% ESCENARIO 3
% ---------------------------------------------------------------------
\subsection{Escenario 3: Sistema de trámites municipales}

\textbf{Actores y responsabilidades.}

{\small
\setlength{\tabcolsep}{4pt}
\renewcommand{\arraystretch}{1.3}
\begin{longtable}{
|>{\raggedright\arraybackslash}p{0.22\textwidth}
|>{\raggedright\arraybackslash}p{0.16\textwidth}
|>{\raggedright\arraybackslash}p{0.50\textwidth}|}
\hline
\textbf{Actor} & \textbf{Tipo} & \textbf{Responsabilidad en el proceso de IR} \\
\hline
\endfirsthead
\hline
\textbf{Actor} & \textbf{Tipo} & \textbf{Responsabilidad en el proceso de IR} \\
\hline
\endhead
Administrador del contrato (GAD) & Cliente / aprobador & Aprueba los entregables, firma las actas y coordina con el municipio. Tiene disponibilidad semanal. \\
\hline
Funcionarios de ventanilla y rentas & Usuarios finales & Describen los flujos de permisos, patentes y certificados, y los cobros. Participan de forma irregular. \\
\hline
Ciudadanos solicitantes & Usuarios finales & Realizan y consultan trámites, y entregan datos personales que deben protegerse. \\
\hline
Personal de tecnología municipal & Stakeholder técnico & Provee acceso a los servidores y conocimiento de la base de datos antigua sin documentación. \\
\hline
Autoridades municipales (ordenanza de tasas) & Stakeholder normativo & La posible ordenanza del cuarto mes puede cambiar tasas y flujos de aprobación. \\
\hline
Entes de contratación y control & Marco contractual & Los pliegos fijan los tres entregables, las fechas, los pagos y las multas por retraso. \\
\hline
Contratista: líder y analista de requisitos & Equipo del proyecto & Elicita, redacta la especificación aprobada en el mes 2, gestiona cambios y mantiene la trazabilidad. \\
\hline
Contratista: desarrollo y pruebas (seis personas en total) & Equipo del proyecto & Implementa e integra los módulos, investiga la base de datos antigua y genera evidencias de pruebas. \\
\hline
\end{longtable}
}

\textbf{Actividades del framework ubicadas en el escenario.}

{\small
\setlength{\tabcolsep}{4pt}
\renewcommand{\arraystretch}{1.3}
\begin{longtable}{
|>{\raggedright\arraybackslash}p{0.20\textwidth}
|>{\raggedright\arraybackslash}p{0.13\textwidth}
|>{\raggedright\arraybackslash}p{0.55\textwidth}|}
\hline
\textbf{Actividad} & \textbf{Tipo} & \textbf{Aplicación en el escenario} \\
\hline
\endfirsthead
\hline
\textbf{Actividad} & \textbf{Tipo} & \textbf{Aplicación en el escenario} \\
\hline
\endhead
Elicitación & Core & Entrevistas con el administrador del contrato, observación del trabajo en ventanilla y rentas, análisis de los pliegos y de la ordenanza vigente, e ingeniería inversa de la base de datos antigua. \\
\hline
Documentación & Core & Especificación de requisitos formal para el mes 2, con trazabilidad, actas de aceptación y registros de pruebas necesarios para justificar los entregables y los pagos. \\
\hline
Negociación y resolución de conflictos & Core & Acordar con el municipio qué cambios (flujo de aprobación, nueva ordenanza) caben dentro del alcance contratado y cuáles requieren una modificación contractual. \\
\hline
Validación & Transversal & Aceptación formal de la especificación en el mes 2, validación periódica de módulos con funcionarios, pruebas de integración y acta de entrega-recepción en el mes 7. \\
\hline
Gestión de requisitos & Transversal & Línea base contractual, control formal de cambios con evaluación de impacto, versionado de la especificación y trazabilidad entre requisitos, pruebas y entregables. \\
\hline
\end{longtable}
}

\textbf{Artefactos del framework.}

{\small
\setlength{\tabcolsep}{4pt}
\renewcommand{\arraystretch}{1.3}
\begin{longtable}{
|>{\raggedright\arraybackslash}p{0.20\textwidth}
|>{\raggedright\arraybackslash}p{0.68\textwidth}|}
\hline
\textbf{Artefacto} & \textbf{Contenido en el escenario} \\
\hline
\endfirsthead
\hline
\textbf{Artefacto} & \textbf{Contenido en el escenario} \\
\hline
\endhead
Metas & Gestionar cerca de 9\,000 trámites anuales; cumplir los tres hitos contractuales sin multas; proteger los datos personales; integrar el sistema con los servidores y la base de datos municipales. \\
\hline
Escenarios & Un ciudadano solicita un permiso de funcionamiento. Un funcionario de rentas cobra una patente con las tasas vigentes. Se aprueba la nueva ordenanza y las tasas cambian. Se emite un certificado de no adeudar consultando la base antigua. \\
\hline
Requisitos orientados a solución & Funcionales: registro, flujo de aprobación, cobro y emisión de documentos. De calidad: seguridad y protección de datos. Restricciones: integración con la base antigua y servidores municipales, y entregables contractuales. \\
\hline
\end{longtable}
}

% ---------------------------------------------------------------------
% ESCENARIO 4
% ---------------------------------------------------------------------
\subsection{Escenario 4: Cooperativa Frutos del Estero}

\textbf{Actores y responsabilidades.}

{\small
\setlength{\tabcolsep}{4pt}
\renewcommand{\arraystretch}{1.3}
\begin{longtable}{
|>{\raggedright\arraybackslash}p{0.22\textwidth}
|>{\raggedright\arraybackslash}p{0.16\textwidth}
|>{\raggedright\arraybackslash}p{0.50\textwidth}|}
\hline
\textbf{Actor} & \textbf{Tipo} & \textbf{Responsabilidad en el proceso de IR} \\
\hline
\endfirsthead
\hline
\textbf{Actor} & \textbf{Tipo} & \textbf{Responsabilidad en el proceso de IR} \\
\hline
\endhead
Gerente de la cooperativa & Cliente / enlace & Disponible casi a diario. Aporta las reglas de negocio y valida funcionalidades, pero no aprueba cambios importantes por sí sola. \\
\hline
Directiva & Decisor & Se reúne mensualmente y aprueba los cambios importantes de alcance. \\
\hline
Socios productores (640) & Usuarios finales & Consultan entregas, pagos, créditos y saldos. Participan principalmente en días de acopio y en asambleas trimestrales. \\
\hline
Personal de acopio & Usuarios internos & Registran peso, humedad y calidad de las entregas, incluso sin conexión. \\
\hline
Compradores exportadores & Stakeholder externo & Sus exigencias de trazabilidad y certificación del cacao pueden cambiar. \\
\hline
Donante & Financiador & Exige informes semestrales de avance, por lo que condiciona los hitos y las evidencias. \\
\hline
Organismo supervisor de la economía popular y solidaria & Marco normativo & Control institucional y cumplimiento del estatuto de la cooperativa. \\
\hline
Equipo de cuatro personas & Equipo del proyecto & Elicita en campo, especifica los cálculos de liquidación, implementa el registro sin conexión y la sincronización, y prepara los informes. \\
\hline
\end{longtable}
}

\textbf{Actividades del framework ubicadas en el escenario.}

{\small
\setlength{\tabcolsep}{4pt}
\renewcommand{\arraystretch}{1.3}
\begin{longtable}{
|>{\raggedright\arraybackslash}p{0.20\textwidth}
|>{\raggedright\arraybackslash}p{0.13\textwidth}
|>{\raggedright\arraybackslash}p{0.55\textwidth}|}
\hline
\textbf{Actividad} & \textbf{Tipo} & \textbf{Aplicación en el escenario} \\
\hline
\endfirsthead
\hline
\textbf{Actividad} & \textbf{Tipo} & \textbf{Aplicación en el escenario} \\
\hline
\endhead
Elicitación & Core & Entrevistas con la gerente y la directiva, observación del proceso actual (tarjetas de papel y hojas de cálculo) en los días de acopio, y consulta a socios en asambleas y recintos. \\
\hline
Documentación & Core & Especificación detallada de reglas de liquidación, créditos y saldos; requisitos de funcionamiento sin conexión y sincronización; documentación de trazabilidad y evidencias para los informes semestrales. \\
\hline
Negociación y resolución de conflictos & Core & Conciliar las prioridades de la gerente, de la directiva y de los socios, y los cambios que piden los compradores con el plazo de doce meses y la aprobación mensual de la directiva. \\
\hline
Validación & Transversal & Validación periódica de incrementos con la gerente, con socios durante los días de acopio y en asambleas trimestrales, y aprobación de la directiva en los cambios importantes. \\
\hline
Gestión de requisitos & Transversal & Línea base para las reglas de liquidación (baja volatilidad), control de cambios de trazabilidad y certificación, y vínculo de cada incremento con los informes del donante. \\
\hline
\end{longtable}
}

\textbf{Artefactos del framework.}

{\small
\setlength{\tabcolsep}{4pt}
\renewcommand{\arraystretch}{1.3}
\begin{longtable}{
|>{\raggedright\arraybackslash}p{0.20\textwidth}
|>{\raggedright\arraybackslash}p{0.68\textwidth}|}
\hline
\textbf{Artefacto} & \textbf{Contenido en el escenario} \\
\hline
\endfirsthead
\hline
\textbf{Artefacto} & \textbf{Contenido en el escenario} \\
\hline
\endhead
Metas & Registrar las entregas de cacao con exactitud; calcular liquidaciones correctas; permitir a los socios consultar entregas, pagos, créditos y saldos; cumplir con el donante y el organismo de control. \\
\hline
Escenarios & Un encargado registra una entrega sin conexión en un recinto y la sincroniza después. Un socio consulta su saldo. La gerente calcula una liquidación. Un comprador exige un nuevo dato de trazabilidad. \\
\hline
Requisitos orientados a solución & Funcionales: registro de entregas, liquidaciones, créditos y consultas. De calidad: exactitud de cálculos, integridad de los datos y protección de datos personales. Restricciones: trabajo sin conexión, sincronización posterior y estatuto de la cooperativa. \\
\hline
\end{longtable}
}

\subsection{Síntesis}

En los cuatro escenarios, el análisis de los actores muestra dónde se concentran las decisiones y qué limita la participación: el jefe de planta y la aseguradora en la Extractora Río Lindo, la propietaria en Sabores de la Abuela Rosa, el administrador del contrato en el GAD Municipal y la directiva junto con el donante en la Cooperativa Frutos del Estero.

Las cinco actividades del framework son las mismas en todos los casos, pero cambia su peso. Cuando predominan la criticidad y las obligaciones formales (Escenarios 1 y 3), destacan la documentación y la gestión de requisitos. Cuando predominan el cambio y la disponibilidad del cliente (Escenario 2), destacan la elicitación continua y la negociación del alcance. En el Escenario 4 la validación debe adaptarse a los días de acopio y a las asambleas. Esta lectura es coherente con los modelos recomendados en C4.


\section{C4. Recomendación fundamentada del modelo de proceso}

Para cada uno de los cuatro escenarios analizados, la selección del modelo de proceso de requisitos se fundamenta en la evaluación de la estabilidad (K1), la disponibilidad de las partes interesadas (K2), la criticidad del sistema (K3), los riesgos principales (K4) y las exigencias normativas o contractuales (K5) \cite{Pohl2025}.

\subsection{Escenario 1: Extractora Río Lindo (Sistema Embebido)}
\begin{itemize}
    \item \textbf{Modelo recomendado:} Modelo Híbrido Controlado (Cascada e Iterativo/Incremental) \cite{GA_Prac_01}.
    \item \textbf{Desglose del modelo híbrido:} Se aplica un enfoque en \textbf{cascada} para los requisitos de seguridad críticos, límites de presión, temperatura y bloqueos definidos por el fabricante de los esterilizadores y exigidos por la aseguradora; y un enfoque \textbf{iterativo e incremental} para el desarrollo de la interfaz de operador, reportes y la integración con los sensores y válvulas aún pendientes de compra \cite{GA_Prac_01}.
    \item \textbf{Razones fundamentadas:}
        \begin{enumerate}
            \item \textit{K1 (Estabilidad):} Las reglas de seguridad y umbrales máximos son estables y fijas del fabricante, permitiendo congelar una línea base temprana \cite{Pohl2025}.
            \item \textit{K5 (Normativa):} La aseguradora exige formalmente un documento de requisitos de seguridad aprobado antes de la puesta en marcha \cite{GA_Prac_01}.
            \item \textit{K3 (Criticidad):} La alta criticidad por la seguridad física de los operadores exige una verificación formal estricta de las funciones de protección \cite{SWEBOK_v4_2024}.
        \end{enumerate}
    \item \textbf{Descarte de los otros modelos:} El modelo ágil puro se descarta porque la alta volatilidad en la interfaz no compensa el riesgo de desatender las reglas críticas de seguridad; el cascada puro resulta inflexible ante la incertidumbre de la selección de sensores en el tercer mes \cite{GA_Prac_01}.
    \item \textbf{Tensión principal:} La incertidumbre en la elección de los proveedores de sensores en el tercer mes frente a la exigencia de una especificación temprana de seguridad \cite{GA_Prac_01}.
\end{itemize}

\subsection{Escenario 2: Sabores de la Abuela Rosa (Aplicación Móvil)}
\begin{itemize}
    \item \textbf{Modelo recomendado:} Modelo Híbrido (Ágil e Iterativo Controlado) \cite{Inayat2015}.
    \item \textbf{Desglose del modelo híbrido:} Se aplica un enfoque \textbf{ágil} basado en un \emph{Product Backlog} vivo para la gestión del catálogo diario de menús, interfaz de usuario y promociones propuestas por la propietaria y su hijo; y un enfoque \textbf{iterativo controlado} para la pasarela de pagos por transferencia y los módulos de seguridad de datos personales conforme a la LOPDP \cite{GA_Prac_01}.
    \item \textbf{Razones fundamentadas:}
        \begin{enumerate}
            \item \textit{K1 (Volatilidad):} El menú diario cambiante y las ideas recurrentes de la dueña generan alta volatilidad que requiere repriorización continua \cite{GA_Prac_01, Pohl2025}.
            \item \textit{K2 (Disponibilidad):} La disponibilidad diaria de Doña Rosa para probar versiones maximiza la validación iterativa frecuente \cite{GA_Prac_01}.
            \item \textit{K4 (Plazo):} El plazo estricto de 14 semanas por las fiestas locales exige entregas incrementales de valor para mitigar el riesgo de descubrimiento tardío \cite{Sommerville2016}.
        \end{enumerate}
    \item \textbf{Descarte de los otros modelos:} El modelo en cascada se descarta por su rigidez ante los cambios diarios del negocio; el iterativo tradicional resulta excesivamente burocrático para un equipo reducido de dos desarrolladores a medio tiempo \cite{GA_Prac_01}.
    \item \textbf{Tensión principal:} La alta agilidad comercial requerida frente a la obligatoriedad estricta de cumplir con la protección legal de datos y conciliación de pagos \cite{GA_Prac_01}.
\end{itemize}

\subsection{Escenario 3: Sistema de Trámites Municipales (Contratación Pública)}
\begin{itemize}
    \item \textbf{Modelo recomendado:} Modelo Híbrido (Iterativo por Hitos Contractuales y Prototipado) \cite{GA_Prac_01}.
    \item \textbf{Desglose del modelo híbrido:} Se aplica un enfoque de \textbf{cascada por fases formales} para cumplir estrictamente con los tres entregables contractuales y de pago (mes 2, 6 y 7); y un enfoque \textbf{iterativo con prototipos} para investigar y resolver la integración técnica con la base de datos municipal antigua y no documentada \cite{GA_Prac_01}.
    \item \textbf{Razones fundamentadas:}
        \begin{enumerate}
            \item \textit{K5 (Contractual):} Los pliegos de contratación pública imponen fechas fijas, entregables aprobados y multas por cada día de retraso \cite{GA_Prac_01}.
            \item \textit{K4 (Riesgo técnico):} La integración con una base de datos heredada sin documentación requiere ciclos de investigación temprana \cite{GA_Prac_01, Pohl2025}.
            \item \textit{K1 (Cambios normativos):} Permite absorber de manera controlada la posible nueva ordenanza de tasas del cuarto mes sin romper la estructura de los entregables \cite{GA_Prac_01}.
        \end{enumerate}
    \item \textbf{Descarte de los otros modelos:} El modelo ágil puro se descarta por las restricciones legales, multas y la inflexibilidad de alcance de la contratación pública; el cascada rígido se descarta porque la ordenanza futura invalidaría una única especificación inicial \cite{GA_Prac_01}.
    \item \textbf{Tensión principal:} La rigidez normativa de los pliegos y multas de contratación pública frente a la incertidumbre técnica de la base de datos heredada \cite{GA_Prac_01}.
\end{itemize}

\subsection{Escenario 4: Cooperativa Frutos del Estero (Plataforma Web Rural)}
\begin{itemize}
    \item \textbf{Modelo recomendado:} Modelo Híbrido (Iterativo-Incremental y Ágil de Campo) \cite{GA_Prac_01}.
    \item \textbf{Desglose del modelo híbrido:} Se aplica un enfoque \textbf{iterativo e incremental} para el núcleo de cálculos de liquidación y créditos (reglas estables del estatuto), combinado con ciclos \textbf{ágiles y de prototipado offline} para el registro de entregas en los centros de acopio rurales con conectividad intermitente \cite{GA_Prac_01}.
    \item \textbf{Razones fundamentadas:}
        \begin{enumerate}
            \item \textit{K2 (Disponibilidad mixta):} Permite sincronizar las revisiones formales con la directiva mensual y la validación de los socios en asambleas trimestrales y días de acopio \cite{GA_Prac_01}.
            \item \textit{K3 (Criticidad técnica):} Facilita comprobar de manera anticipada el funcionamiento sin conexión y la posterior sincronización de datos \cite{GA_Prac_01}.
            \item \textit{K5 (Informes):} La generación de incrementos funcionales semestrales respalda directamente los informes exigidos por el donante \cite{GA_Prac_01}.
        \end{enumerate}
    \item \textbf{Descarte de los otros modelos:} El modelo ágil se descarta porque la directiva se reúne solo una vez al mes y los socios están dispersos, impidiendo la priorización continua; el cascada no permitiría validar oportunamente el funcionamiento offline \cite{GA_Prac_01}.
    \item \textbf{Tensión principal:} Las exigencias de reportes formales semestrales del donante frente a la baja frecuencia de retroalimentación directa de los productores rurales \cite{GA_Prac_01}.
\end{itemize}

\section{C5. Análisis de riesgos del modelo elegido y mitigación}

\begin{enumerate}
    \item \textbf{Riesgo 1 en Escenario 1 (Híbrido controlado: Cascada + Iterativo/Incremental):} Obsolescencia de los requisitos ante cambios en los sensores seleccionados en el tercer mes \cite{GA_Prac_01}.
        \begin{itemize}
            \item \textit{Impacto:} Afecta el diseño del hardware y genera retrabajo en la integración del sistema embebido.
            \item \textit{Mitigación:} Establecer una capa de abstracción de hardware (HAI) que aísle el software de control de la marca específica de sensores.
        \end{itemize}
    \item \textbf{Riesgo 2 en Escenario 1 (Híbrido controlado: Cascada + Iterativo/Incremental):} Dificultad de validación continua por la disponibilidad baja a media del jefe de planta (dos mañanas por semana) y turnos rotativos de operadores \cite{GA_Prac_01}.
        \begin{itemize}
            \item \textit{Impacto:} Retrasos en la aprobación formal de los flujos de control y detección tardía de malentendidos operativos.
            \item \textit{Mitigación:} Programar revisiones estructuradas fijas al cierre de cada incremento y validar mediante prototipos de interfaz antes de la codificación final.
        \end{itemize}

    \item \textbf{Riesgo 1 en Escenario 2 (Híbrido: Ágil + Iterativo controlado):} Omisión de Requisitos No Funcionales (RNF) de seguridad y cumplimiento de la LOPDP en la aplicación móvil \cite{GA_Prac_01}.
        \begin{itemize}
            \item \textit{Impacto:} Incumplimiento de la normativa legal ecuatoriana de protección de datos personales de los clientes.
            \item \textit{Mitigación:} Incorporar historias de usuario de seguridad y privacidad como criterios de aceptación obligatorios desde el primer sprint \cite{Pohl2025}.
        \end{itemize}
    \item \textbf{Riesgo 2 en Escenario 2 (Híbrido: Ágil + Iterativo controlado):} Expansión descontrolada del alcance (\emph{scope creep}) debido a las ideas recurrentes de la propietaria y su hijo en cada reunión \cite{GA_Prac_01}.
        \begin{itemize}
            \item \textit{Impacto:} Agotamiento prematuro del presupuesto cerrado y retraso en la entrega para las fiestas locales.
            \item \textit{Mitigación:} Gestionar rigurosamente el \emph{Product Backlog} aplicando priorización estricta (técnica MoSCoW) para congelar el alcance de cada incremento.
        \end{itemize}

    \item \textbf{Riesgo 1 en Escenario 3 (Híbrido: Cascada por hitos + Iterativo con prototipado):} Desviación del alcance acordado en el entregable del mes 2 debido a modificaciones en el flujo de aprobación del GAD \cite{GA_Prac_01}.
        \begin{itemize}
            \item \textit{Impacto:} Riesgo de retrasos y aplicación de multas contractuales estipuladas en los pliegos.
            \item \textit{Mitigación:} Gestionar los cambios mediante un comité formal de control de cambios (CCB) documentado en actas de iteración \cite{Pohl2025}.
        \end{itemize}
    \item \textbf{Riesgo 2 en Escenario 3 (Híbrido: Cascada por hitos + Iterativo con prototipado):} Incompatibilidad o errores de lectura en la base de datos municipal antigua y no documentada \cite{GA_Prac_01}.
        \begin{itemize}
            \item \textit{Impacto:} Fallas críticas en la consulta de datos catastrales durante la emisión de trámites.
            \item \textit{Mitigación:} Dedicar la primera iteración del contrato a realizar ingeniería inversa y pruebas de concepto (PoC) sobre la estructura de la base de datos heredada.
        \end{itemize}

    \item \textbf{Riesgo 1 en Escenario 4 (Híbrido: Iterativo-Incremental + Ágil de campo):} Fallas en la sincronización de registros de cacao debido a la inestabilidad de internet en el centro de acopio de Mocache \cite{GA_Prac_01}.
        \begin{itemize}
            \item \textit{Impacto:} Pérdida temporal de datos de entrega y descontento en los pagos a los productores.
            \item \textit{Mitigación:} Priorizar pruebas tempranas de almacenamiento local offline (IndexedDB/SQLite) y colas de sincronización robustas en los primeros incrementos.
        \end{itemize}
    \item \textbf{Riesgo 2 en Escenario 4 (Híbrido: Iterativo-Incremental + Ágil de campo):} Desalineación entre los requerimientos de la directiva (reunión mensual) y las necesidades reales de los socios dispersos (asambleas trimestrales) \cite{GA_Prac_01}.
        \begin{itemize}
            \item \textit{Impacto:} Rechazo de la plataforma por parte de los usuarios finales por no ajustarse a su operación en recintos rurales.
            \item \textit{Mitigación:} Nombrar comités de enlace durante los días de acopio para validar incrementos funcionales periódicos que respalden los informes semestrales del donante.
        \end{itemize}
\end{enumerate}

\begin{thebibliography}{9}
\bibitem{Pohl2025} Klaus Pohl. \emph{Requirements Engineering: Fundamentals, Principles, and Techniques}. Springer, 2nd edition, 2025.
\bibitem{IREB2026} Bühne, S., Glinz, M., van Loenhoud, H. y Staal, S. (2026). \emph{CPRE Foundation Level Syllabus, versión 3.3.0}. International Requirements Engineering Board. \url{https://www.ireb.org/en/downloads/}
\bibitem{ISO29148} ISO/IEC/IEEE. (2018). \emph{ISO/IEC/IEEE 29148:2018, Systems and software engineering — Life cycle processes — Requirements engineering}. \url{https://webstore.iec.ch/en/publication/64315}
\bibitem{SWEBOK_v4_2024} Hironori Washizaki, editor. \emph{Guide to the Software Engineering Body of Knowledge (SWEBOK Guide)}, Version 4.0. IEEE Computer Society, 2024.
\bibitem{Inayat2015} Irum Inayat et al. A systematic literature review on agile requirements engineering practices and challenges. \emph{Computers in Human Behavior}, 51:915-929, 2015.
\bibitem{Sommerville2016} Ian Sommerville. \emph{Software Engineering}. Pearson, 10th edition, 2016.
\bibitem{GA_Prac_01} Universidad Técnica Estatal de Quevedo (UTEQ). \emph{Material complementario GA-PRAC-01: Proceso de requisitos}, 2026.
\end{thebibliography}


\end{document}
